\documentclass[10pt,a4paper]{amsart}
\usepackage[margin=19mm]{geometry}
\usepackage{amsmath,amssymb,mathtools}
\usepackage[scr=rsfs,cal=euler]{mathalfa}
\usepackage{xfrac}
\usepackage[unicode,hidelinks,bookmarks=false]{hyperref}
\newcommand{\ie}{i.e.\ }
\newcommand{\cf}{cf.\ }
\newcommand{\resp}{respectively\ }
\newcommand{\Subsection}{\subsection}
\providecommand{\nobreakdash}{\nobreak\hbox{-}}
\setlength{\emergencystretch}{2em}
\multlinegap0pt
\theoremstyle{plain}
\newtheorem{theorem}{Theorem}
\newtheorem{lemma}[theorem]{Lemma}
\newtheorem{corollary}[theorem]{Corollary}
\newtheorem{proposition}[theorem]{Proposition}
\newtheorem{claim}[theorem]{Claim}
\theoremstyle{definition}
\newtheorem{definition}[theorem]{Definition}
\theoremstyle{remark}
\newtheorem{remark}[theorem]{Remark}
\newtheorem{assumption}{Assumption}
\title[An algebraic approach to product-form ]{An algebraic approach to product-form stationary distributions for some reaction networks: the kernel over the field of fractions}
\author{Beatriz Pascual-Escudero; Linard Hoessly}
\date{Source: arXiv:2012.03227v2 (CC BY 4.0)}
\def\Ga{\Gamma}
\def\R{\mathbb{R}}
\def\cR{\mathcal{R}}
\def\k{\kappa}
\def\la{\lambda}
\def\p{^{\prime}}
\begin{document}
\maketitle
\noindent\emph{Source and licence.} An algebraic approach to product-form stationary distributions for some reaction networks. Version arXiv:2012.03227v2 (CC BY 4.0), distributed under the Creative Commons Attribution 4.0 International licence, \url{http://creativecommons.org/licenses/by/4.0/}, as recorded in the arXiv licence metadata for this exact version and shown on the abstract page \url{https://arxiv.org/abs/2012.03227v2}. Full licence notice: licenses/arxiv-2012.03227-CC-BY-4.0.txt.\par\smallskip

\noindent\textit{Editorial scope.} Counted unit: the article's Corollary with source label \texttt{clm2} (source0.tex lines 502--504), which gives the kernel of the reduced matrix, under the article's linearity assumption, as a unique vector of homogeneous polynomials up to multiplication by a constant; delivered together with the article's complete original proof of it (source0.tex lines 506--538, one \texttt{proof} environment of 33 lines, adjacent to the statement). No mathematics is written, repaired or re-derived: statement and proof are the article's own text line for line, in the article's own notation (the reaction network, the reduced matrix, the field of fractions, the kernel and the homogeneity), and no retained line is substituted, re-flowed or omitted. Required context retained uncounted, in source order: the statement of the article's linearity assumption (lines 459--461) and the statement of the article's rank lemma (line 494, a complete one-line lemma in the article), both of which the counted statement and proof invoke. The proofs of those context results are not delivered. Nothing else of the article is retained: its main theorems, its examples and its remaining lemmas are omitted. No citation command occurs in any retained line, so the article's bibliography (an inline bibliography that the article ships) is not delivered and the wrapper carries no bibliography entry. Every cross-reference inside the retained spans resolves inside the delivered document. Editorial formatting only: an AMS \texttt{amsart} preamble that restates the environments and the macros the retained lines use, namely the article's declarations of Theorem, Lemma, Corollary, Proposition, Claim, Definition, Remark and Assumption, and the macros for the special letters of the retained lines. The article's own source declares the class \texttt{amsart} as well; no class or style file is shipped with this package. Numbering differs from the article, because the article declares its Lemma environment as the master counter, resets it in each section, and only a subset is retained: the retained assumption prints as Assumption 1, the retained rank lemma as Lemma 1 and the counted corollary as Corollary 2. No picture environment and no graphical inclusion occurs in any retained line, and the package delivers no image asset. One bundled source file, \texttt{sources/105523/source0.tex}, accompanies the delivered item as the exact source of every retained span. Every retained span is recorded in \texttt{delivery-evidence.json} with method \texttt{exact\_tex}.
\begin{assumption}\label{ass1}
The transition intensity functions $\la_{\nu \to \nu\p}(\cdot)$ are linear polynomials in the reaction rate parameters $\k_{\nu \to \nu\p}$ for all reactions $\nu \to \nu\p\in\cR$. 
\end{assumption}

 \begin{lemma}\label{clm}  Consider $A(\k)$ as a matrix with entries in the field of fractions $\R(\k_i|{\nu_i\to\nu_i\p}\in\cR)$. Under Assumption \ref{ass1}, $\mathrm{rank}(A(\k))=m-1$\end{lemma}

\begin{corollary}\label{clm2}
Under Assumption \ref{ass1}, the map $h_\Ga$ can be given as a unique vector $(F_1,\cdots , F_m)$, up to multiplication by a constant in $\R^*$, whose entries are homogeneous polynomials in $\R[\k_i|{\nu_i\to\nu_i\p}\in\cR]$ of the same degree, with positive coefficients and no non-constant common divisor.  Hence the $F_i$ never vanish on $\R\Pr_{>0}^\cR$. 
\end{corollary}

\begin{proof}
For the ease of the reader we complete the computation of $h_\Ga$ by recalling the computation of the kernel of $A(\k)$ (which is one-dimensional by Lemma \ref{clm}) via Gaussian elimination over the field of fractions $\R(\k_i|{\nu_i\to\nu_i\p}\in\cR)$. This gives a constructive proof for $h_\Ga$ and $(F_1,\cdots,F_m)$. We can proceed as follows:
$$(0)\quad \left[
    	\begin{array}{cc|cc}
    	A(\k) \\
    	\hline
    	I_m \\	
    	\end{array}
    	\right]\xRightarrow{(1)\&(2)}
	\left[
    	\begin{array}{cc|cc}
    	B \\
    	\hline
    	C \\	
    	\end{array}
    	\right]\quad (3),	$$


\begin{enumerate}
\setcounter{enumi}{-1}
\item We start by creating the row-augmented matrix, where $I_m$ is the $m\times m$ identity matrix.
\item We divide each j-column by $-a_{j,j}=\sum_{\nu_i\to\nu_i\p\in\cR}\la_i(x_j)$. Hence we get $-1$-entries in the diagonal of the former matrix $A(\k)$. 
\item We apply elementary column operations to get the upper matrix in column echelon form. Note that we only add columns multiplied by elements in $\R(\k_i|{\nu_i\to\nu_i\p}\in\cR)$ with all coefficients positive.
\item We get the matrices $B,C$, where $B$ is in column echelon form of degree $m-1$. Furthermore  matrix $C$ has the form $$C=\begin{pmatrix}
c_{1,1}\,\cdots\, c_{1,m-1}\, v_1\\
\vdots\,\ddots\,\quad\,\vdots\,\,\,\,\,\,\,\,\,\,\,\,\,\,\,\,\vdots\\
c_{m,1}\cdots\, c_{m,m-1}\, v_m
\end{pmatrix}$$
	such that the last column of $C$, i.e. $(v_1,\cdots, v_m)$, is a non-trivial vector of the kernel.
\end{enumerate}
Hence $(v_1,\cdots, v_m)$ is a basis of the kernel over $\R(\k_i|{\nu_i\to\nu_i\p}\in\cR)$, and we can choose the corresponding column vector as homogeneous polynomials. If the greatest common divisor(GCD) in $\R[\k_i|{\nu_i\to\nu_i\p}\in\cR]$ is nontrivial we divide all polynomials of the column vector by the GCD. Denote the corresponding vector by $w$ with entries in $\R[\k_i|{\nu_i\to\nu_i\p}\in\cR]$, which is unique up to multiplication by a constant in $\R^*$.
 Then the vector $w$ has the form $w=(F_1,\cdots ,F_m)$, and by construction the $F_1,\cdots ,F_m$ are homogeneous polynomials of the same degree with positive coefficients by (2). As all the coefficients of the polynomials are positive, it is clear that $F_1,\cdots, F_m$ never vanish on $\R\Pr_{>0}^\cR$.
\end{proof}

\end{document}
